\subsection{GROUPS OPERATING FREELY ON A TOPOLOGICAL SPACE}

DEFINITION 2. Let G be a group operating on a set X. G is said to operate freely on X if the stabilizer of every element of X is \(\{e\}\) , in other words if the relations \(s.x = x, s \in G, x \in X\) imply s = e.

Example. Let G be a group and let H be a subgroup of G. Then H operates freely by (left or right) translations on G.

Let G be a group operating freely on a set X, let R be the equivalence relation defined by G on X, and let \(C \subset X \times X\) be the graph of R. If \((x, y) \in C\) , then there exists \(s \in G\) such that s.x = y; and s is unique, since \(s.x = s'.x\) implies \(s'^{-1}s.x = x\) , and therefore \(s'^{-1}s = e\) (as G operates freely). If we make correspond to \((x, y) \in C\) the unique \(s \in G\) such that s.x = y, we define a mapping \(\varphi: C \to G\) , which we shall call the canonical mapping of C into G. With this notation:

PROPOSITION 6. Let G be a topological group operating continuously on a topological space X, and suppose that G operates freely on X. Then G operates properly on X if and only if the following condition is satisfied:

(FP) The graph C of the equivalence relation defined by G is closed in \(X \times X\) , and the canonical mapping \(\varphi: C \rightarrow G\) is continuous.

The set C is the image of the mapping \(\theta: (s, x) \to (x, s.x)\) of \(\mathbf{G} \times \mathbf{X}\) into \(\mathbf{X} \times \mathbf{X}\) . We know (Chapter I, § 10, no. 1, Proposition 2) that \(\theta\) is proper if and only if C is closed in \(\mathbf{X} \times \mathbf{X}\) and (if \(\theta'\) denotes the mapping \(\theta\) considered as a mapping of \(\mathbf{G} \times \mathbf{X}\) into C) \(\theta'\) is a homeomorphism. Now the hypothesis implies that \(\theta'\) is bijective and that its inverse is the mapping \((x, y) \to (\varphi(x, y), x)\) . Hence \(\theta'\) is a homeomorphism if and only if \(\varphi\) is continuous.

